% thm-shape: limit-shapes.tex:106-119 (label thm:shape)
\paragraph{Paper excerpt.}
\begin{verbatim}
\begin{theorem}[Limit shape]\label{thm:shape}
Let $d\geq2$ and $0<\drift< \frac1d$. Then, almost surely, the centrally excited random walk has the following properties.
\begin{enumerate}[label=\textup{(\roman*)}]
\item \underline{\emph{Shape}}: for every $0<\eta<1$ and all sufficiently large~$n$,
\begin{equation*}
\{x\in\Z^d:|x|<(1-\eta)r_n\}\subset A_n\subset\{x\in\Z^d:|x|<(1+\eta)r_n\}\, .
\end{equation*}
\item \underline{\emph{Local times}}: as $n\to\infty$,
\begin{equation*}
\frac1{r_n}\max_{x\in\Z^d}\bigl|\ell_n(x)-2d\drift(r_n-|x|)_+\bigr|\longrightarrow0\, .
\end{equation*}
\item \underline{\emph{Recurrence}}: the walk visits every site of~$\Z^d$ infinitely often.
\end{enumerate}
\end{theorem}
\end{verbatim}
\paragraph{Lean statement.}
\begin{verbatim}
theorem CERW.Frozen.limit_shape {d : ℕ} (hd : 2 ≤ d)
    {ε : ℝ} (hε : 0 < ε) (hεd : ε < 1 / (d : ℝ))
    {Ω : Type u} [MeasurableSpace Ω] (μ : Measure Ω) [IsProbabilityMeasure μ]
    (X : ℕ → Ω → Site d) (hX : CERW.IsCERW μ ε X) :
    let ωd : ℝ := (volume (Metric.ball (0 : EuclideanSpace ℝ (Fin d)) 1)).toReal
    let r : ℕ → ℝ := fun n => ((d + 1) * n / (2 * d * ε * ωd)) ^ ((1 : ℝ) / (d + 1))
    ∀ᵐ ω ∂μ,
      (∀ η : ℝ, 0 < η → η < 1 → ∀ᶠ n : ℕ in atTop,
        {x : Site d | euclidNorm x < (1 - η) * r n} ⊆ ↑(CERW.departureRange (X · ω) n) ∧
        (↑(CERW.departureRange (X · ω) n) : Set (Site d)) ⊆
          {x | euclidNorm x < (1 + η) * r n}) ∧
      (∀ η : ℝ, 0 < η → ∀ᶠ n : ℕ in atTop, ∀ x : Site d,
        |(CERW.localTime (X · ω) n x : ℝ) - 2 * d * ε * max (r n - euclidNorm x) 0|
          ≤ η * r n) ∧
      (∀ x : Site d, ∃ᶠ j in atTop, X j ω = x)
--
\end{verbatim}
\paragraph{Transcription.}
The walk is any process \texttt{X} on any probability space $(\Omega,\mu)$ with \texttt{CERW.IsCERW μ ε X} ($\varepsilon=\kappa$); $d\ge2$ is \texttt{hd : 2 ≤ d} and $0<\kappa<1/d$ is \texttt{hε, hεd}. $\omega_d$ is \texttt{(volume (ball 0 1)).toReal} and $r_n$ is the \texttt{let} \texttt{r}, with the natural-number casts taken in $\mathbb R$. $\ell_n$ is \texttt{CERW.localTime} and $A_n$ is \texttt{CERW.departureRange} (coerced to a \texttt{Set}), so $A_n=\{X_0,\dots,X_{n-1}\}$ as in the paper. ``Almost surely'' is a single \texttt{∀ᵐ ω ∂μ} over the conjunction of (i)--(iii), with the quantifier over $\eta$ inside it. (i) is ``for every $0<\eta<1$, eventually'' with the two inclusions. (ii) ``$\frac1{r_n}\max_x|\cdot|\to0$'' is read as $\forall\eta>0$, eventually in $n$, $\forall x$, $|\ell_n(x)-2d\kappa\max(r_n-|x|,0)|\le\eta r_n$ (for $|x|\ge r_n$ outside the range both $\ell_n(x)$
